\chapter{1993年全国卷（文）}
\section{单选题}

\begin{problem}
如果双曲线的实半轴长为 \(2\)，焦距为 \(6\)，那么该双曲线的离心率为（\quad）

\choices
    {\(\frac{\sqrt{3}}{2}\)}
    {\(\frac{\sqrt{6}}{2}\)}
    {\(\frac{3}{2}\)}
    {\(2\)}
\end{problem}

\begin{answer}
C
\end{answer}

\begin{solution}
双曲线的实半轴长为 \(a=2\)，焦距为 \(2c=6\)，所以 \(c=3\). 因此离心率为 \(e=\frac{c}{a}=\frac{3}{2}\)，故选 C.
\end{solution}

\begin{problem}
函数 \(y = \frac{1 - \tan^2 2x}{1 + \tan^2 2x} \) 的最小正周期是（\quad）

\choices
    {\(\frac{\pi}{4}\)}
    {\(\frac{\pi}{2}\)}
    {\(\pi\)}
    {\(2\pi\)}
\end{problem}

\begin{answer}
B
\end{answer}

\begin{solution}
在函数的定义域内，
\(\displaystyle \frac{1-\tan^2 2x}{1+\tan^2 2x}=\cos 4x\). 原式的定义域排除了 \(x=\frac{\pi}{4}+k\frac{\pi}{2}\)，这个排除点集的最小正周期也是 \(\frac{\pi}{2}\)，而 \(\cos 4x\) 的最小正周期为 \(\frac{\pi}{2}\)，故选 B.
\end{solution}

\begin{problem}
当圆锥的侧面积和底面积的比值是 \(\sqrt{2}\) 时，圆锥的轴截面顶角是（\quad）

\choices
    {\(45^{\circ}\)}
    {\(60^{\circ}\)}
    {\(90^{\circ}\)}
    {\(120^{\circ}\)}
\end{problem}

\begin{answer}
C
\end{answer}

\begin{solution}
设圆锥的底面半径为 \(r\)，母线长为 \(l\). 由侧面积与底面积之比得
\(\displaystyle \frac{\pi rl}{\pi r^2}=\frac{l}{r}=\sqrt{2}\)，所以 \(\sin\frac{\theta}{2}=\frac{r}{l}=\frac{1}{\sqrt{2}}\)，其中 \(\theta\) 是轴截面顶角. 由于 \(0<\frac{\theta}{2}<\frac{\pi}{2}\)，有 \(\frac{\theta}{2}=45^\circ\)，从而 \(\theta=90^\circ\)，故选 C.
\end{solution}

\begin{problem}
当 \(z = \frac{1 + \i}{\sqrt{2}}\) 时，\(z^{100} + z^{50} + 1\) 的值等于（\quad）

\choices
    {\(1\)}
    {\(-1\)}
    {\(\i\)}
    {\(-\i\)}
\end{problem}

\begin{answer}
C
\end{answer}

\begin{solution}
\(\displaystyle z=\frac{1+\i}{\sqrt{2}}=\cos\frac{\pi}{4}+\i\sin\frac{\pi}{4}\)，所以
\(z^{50}=\cos\frac{25\pi}{2}+\i\sin\frac{25\pi}{2}=\i\)，且 \(z^{100}=\cos25\pi+\i\sin25\pi=-1\). 因此 \(z^{100}+z^{50}+1=-1+\i+1=\i\)，故选 C.
\end{solution}

\begin{problem}
若正棱锥的底面边长与侧棱长相等，则该棱锥一定不是（\quad）

\choices
    {三棱锥}
    {四棱锥}
    {五棱锥}
    {六棱锥}
\end{problem}

\begin{answer}
D
\end{answer}

\begin{solution}
设正棱锥的底面为正 \(n\) 边形，底面边长和侧棱长均为 \(s\)，底面中心到顶点的距离为底面外接圆半径 \(R=\frac{s}{2\sin(\pi/n)}\). 若棱锥顶点到底面中心的高为 \(h\)，则由直角三角形得 \(s^2=R^2+h^2\)，所以必须有 \(R<s\). 当 \(n=3\)，\(4\)，\(5\) 时，\(\sin\frac{\pi}{n}>\frac{1}{2}\)，可以取正的 \(h\)；当 \(n=6\) 时，\(R=s\)，只能有 \(h=0\)，不能构成棱锥. 因此该棱锥一定不是六棱锥，故选 D.
\end{solution}

\begin{problem}
在直角三角形中两锐角为 \(A\) 和 \(B\)，则 \(\sin A\sin B\)（\quad）

\choices
    {有最大值 \(\frac{1}{2}\) 和最小值 \(0\)}
    {有最大值 \(\frac{1}{2}\)，但无最小值}
    {既无最大值，也无最小值}
    {有最大值 \(1\)，但无最小值}
\end{problem}

\begin{answer}
B
\end{answer}

\begin{solution}
直角三角形的两锐角满足 \(A+B=90^\circ\)，于是
\(\sin A\sin B=\sin A\cos A=\frac{1}{2}\sin 2A\). 当 \(0<A<90^\circ\) 时，\(0<2A<180^\circ\)，所以该式的最大值为 \(\frac{1}{2}\)，在 \(A=45^\circ\) 时取得；当 \(A\) 趋近于 \(0^\circ\) 或 \(90^\circ\) 时，该式趋近于 \(0\)，但锐角不能取这两个端点，因而无最小值，故选 B.
\end{solution}

\begin{problem}
在各项均为正数的等比数列 \(\{a_{n}\}\) 中，若 \(a_{5}a_{6} = 9\)，则 \(\log_{3} a_{1} + \log_{3} a_{2} + \cdots + \log_{3} a_{10}\) 的值为（\quad）

\choices
    {\(12\)}
    {\(10\)}
    {\(8\)}
    {\(2 + \log_{3} 5\)}
\end{problem}

\begin{answer}
B
\end{answer}

\begin{solution}
设等比数列的首项为 \(a_1\)，公比为 \(q>0\). 则
\(a_5a_6=a_1^2q^9=9\). 另一方面，
\(\displaystyle \prod_{k=1}^{10}a_k=a_1^{10}q^{0+1+\cdots+9}=a_1^{10}q^{45}=(a_5a_6)^5=9^5=3^{10}\). 因此
\(\displaystyle \log_3a_1+\log_3a_2+\cdots+\log_3a_{10}=\log_3\left(\prod_{k=1}^{10}a_k\right)=10\)，故选 B.
\end{solution}

\begin{problem}
\(F(x) = \left(1 + \frac{2}{2^x - 1}\right)f(x)\)（\(x\neq 0\)）是偶函数，且 \(f(x)\) 不恒等于零，则 \(f(x)\)（\quad）

\choices
    {是奇函数}
    {是偶函数}
    {可能是奇函数也可能是偶函数}
    {不是奇函数也不是偶函数}
\end{problem}

\begin{answer}
A
\end{answer}

\begin{solution}
令 \(\displaystyle g(x)=1+\frac{2}{2^x-1}=\frac{2^x+1}{2^x-1}\). 当 \(x\ne0\) 时，
\(\displaystyle g(-x)=\frac{2^{-x}+1}{2^{-x}-1}=-\frac{2^x+1}{2^x-1}=-g(x)\). 由 \(F(x)=g(x)f(x)\) 为偶函数，得
\(g(x)f(x)=F(x)=F(-x)=g(-x)f(-x)=-g(x)f(-x)\). 又 \(g(x)\ne0\)，所以 \(f(x)=-f(-x)\)，即 \(f(x)\) 是奇函数，故选 A.
\end{solution}

\begin{problem}
设直线 \(2x - y - \sqrt{3} = 0\) 与 \(y\) 轴的交点为 \(P\)，点 \(P\) 把圆 \((x + 1)^2 + y^2 = 25\) 的直径分为两段，则其长度之比为（\quad）

\choices
    {\(7:3\) 或 \(3:7\)}
    {\(7:4\) 或 \(4:7\)}
    {\(7:5\) 或 \(5:7\)}
    {\(7:6\) 或 \(6:7\)}
\end{problem}

\begin{answer}
A
\end{answer}

\begin{solution}
令直线与 \(y\) 轴的交点为 \(P\)，代入 \(x=0\) 得 \(P=(0,-\sqrt{3})\). 圆心为 \(O=(-1,0)\)，半径为 \(5\)，且 \(OP=\sqrt{1+3}=2\). 过 \(P\) 的任一直径的两个端点到圆心的距离均为 \(5\)，而 \(P\) 位于圆内，所以它把该直径分成长度 \(5+2=7\) 和 \(5-2=3\) 的两段，长度之比为 \(7:3\) 或 \(3:7\)，故选 A.
\end{solution}

\begin{problem}
若 \(a\)，\(b\) 是任意实数，且 \(a > b\)，则（\quad）

\choices
    {\(a^2 > b^2\)}
    {\(\frac{b}{a} < 1\)}
    {\(\lg (a - b) > 0\)}
    {\(\left(\frac{1}{2}\right)^a < \left(\frac{1}{2}\right)^b\)}
\end{problem}

\begin{answer}
D
\end{answer}

\begin{solution}
由 \(a>b\) 不能推出 \(a^2>b^2\)，例如取 \(a=-1\)，\(b=-2\)，则 \(a^2=1<4=b^2\)；同样取 \(a=-1\)，\(b=-2\)，则 \(\frac{b}{a}=2>1\)，所以选项 B 也不一定成立；取 \(a=2\)，\(b=1\)，则 \(\lg(a-b)=\lg1=0\)，所以选项 C 不成立. 因为底数 \(\frac{1}{2}\) 在 \((0,+\infty)\) 上是减函数，且 \(a>b\)，所以 \(\left(\frac{1}{2}\right)^a<\left(\frac{1}{2}\right)^b\)，故选 D.
\end{solution}

\begin{problem}
已知集合 \(E = \{\theta \mid \cos \theta < \sin \theta,\ 0 \leqslant \theta \leqslant 2\pi\}\)，\(F = \{\theta \mid \tan \theta < \sin \theta\}\)，那么 \(E \cap F\) 为区间（\quad）

\choices
    {\(\left(\frac{\pi}{2},\pi\right)\)}
    {\(\left(\frac{\pi}{4},\frac{3\pi}{4}\right)\)}
    {\(\left(\pi, \frac{3\pi}{2}\right)\)}
    {\(\left(\frac{3\pi}{4}, \frac{5\pi}{4}\right)\)}
\end{problem}

\begin{answer}
A
\end{answer}

\begin{solution}
由 \(\cos\theta<\sin\theta\) 得 \(\sin\left(\theta-\frac{\pi}{4}\right)>0\)，结合 \(0\leqslant\theta\leqslant2\pi\)，可得 \(E=\left(\frac{\pi}{4},\frac{5\pi}{4}\right)\). 对 \(F\)，在正切有定义时，
\(\tan\theta-\sin\theta=\frac{\sin\theta(1-\cos\theta)}{\cos\theta}\). 在 \(0\leqslant\theta\leqslant2\pi\) 内，该式在 \(\left(\frac{\pi}{2},\pi\right)\) 和 \(\left(\frac{3\pi}{2},2\pi\right)\) 内为负，在其余正切有定义的区间内不为负，因此 \(F\cap[0,2\pi]=\left(\frac{\pi}{2},\pi\right)\cup\left(\frac{3\pi}{2},2\pi\right)\). 于是 \(E\cap F=\left(\frac{\pi}{2},\pi\right)\)，故选 A.
\end{solution}

\begin{problem}
一动圆与两圆：\(x^2 + y^2 = 1\) 和 \(x^2 + y^2 - 8x + 12 = 0\) 都外切，则动圆圆心的轨迹为（\quad）

\choices
    {抛物线}
    {圆}
    {双曲线的一支}
    {椭圆}
\end{problem}

\begin{answer}
C
\end{answer}

\begin{solution}
设动圆圆心为 \(X\)，半径为 \(r\). 两个已知圆的圆心分别为 \(O_1=(0,0)\)、\(O_2=(4,0)\)，半径分别为 \(1\)、\(2\). 由外切条件得 \(XO_1=r+1\)，\(XO_2=r+2\)，所以 \(XO_2-XO_1=1\). 圆心 \(O_1\)，\(O_2\) 的距离为 \(4\)，动点到两定点距离之差的绝对值为定值 \(1\)，且小于两焦点间距离，故轨迹为双曲线的一支，选 C.
\end{solution}

\begin{problem}
若直线 \(ax + by + c = 0\) 在第一，二，三象限，则（\quad）

\choices
    {\(ab > 0\)，\(bc > 0\)}
    {\(ab > 0\)，\(bc < 0\)}
    {\(ab < 0\)，\(bc > 0\)}
    {\(ab < 0\)，\(bc < 0\)}
\end{problem}

\begin{answer}
D
\end{answer}

\begin{solution}
直线不可能是竖直线或水平线，因此可写成 \(y=mx+n\). 要同时经过第一、三象限，必须有 \(m>0\)；要在第二象限取得点，还必须有 \(n>0\). 由 \(m=-\frac{a}{b}>0\) 得 \(ab<0\)，由 \(n=-\frac{c}{b}>0\) 得 \(bc<0\)，故选 D.
\end{solution}

\begin{problem}
如果圆柱轴截面的周长 \(l\) 为定值，那么圆柱体积的最大值是（\quad）

\choices
    {\(\left(\frac{l}{6}\right)^3\pi\)}
    {\(\left(\frac{l}{3}\right)^3\pi\)}
    {\(\left(\frac{l}{4}\right)^3\pi\)}
    {\(\frac{1}{4}\left(\frac{l}{4}\right)^3\pi\)}
\end{problem}

\begin{answer}
A
\end{answer}

\begin{solution}
设圆柱的底面半径为 \(r\)，高为 \(h\). 轴截面是长为 \(h\)、宽为 \(2r\) 的矩形，所以 \(2(h+2r)=l\)，即 \(h+2r=\frac{l}{2}\). 令 \(u=2r\)，则 \(r^2h=\frac{u^2h}{4}\)，而 \(\frac{u}{2}+\frac{u}{2}+h=\frac{l}{2}\). 由均值不等式，\(\displaystyle r^2h=\frac{u}{2}\cdot\frac{u}{2}\cdot h\leqslant\left(\frac{l/2}{3}\right)^3=\left(\frac{l}{6}\right)^3\)，等号当且仅当 \(r=h=\frac{l}{6}\) 时成立. 因此圆柱体积的最大值为 \(\left(\frac{l}{6}\right)^3\pi\)，故选 A.
\end{solution}

\begin{problem}
由 \((\sqrt{3} x + \sqrt[3]{2})^{100}\) 展开所得的 \(x\) 的多项式中，系数为有理数的共有（\quad）

\choices
    {\(50\) 项}
    {\(17\) 项}
    {\(16\) 项}
    {\(15\) 项}
\end{problem}

\begin{answer}
B
\end{answer}

\begin{solution}
展开式的通项为 \(T_{k+1}=\mathrm C_{100}^{k}(\sqrt{3}x)^k(\sqrt[3]{2})^{100-k}\)，其中 \(k=0\)，\(1\)，\(\ldots\)，\(100\). 其 \(x^k\) 的系数为有理数，当且仅当 \(k\) 为偶数且 \(100-k\) 是 \(3\) 的倍数；否则分别会含有不可消去的 \(\sqrt{3}\) 因子或 \(\sqrt[3]{2}\) 因子. 于是 \(k\equiv4\pmod{6}\)，在 \(0\leqslant k\leqslant100\) 中为 \(4\)，\(10\)，\(\ldots\)，\(100\)，共有 \(\frac{100-4}{6}+1=17\) 个，故选 B.
\end{solution}

\begin{problem}
设 \(a\)，\(b\)，\(c\) 都是正数，且 \(3a = 4b = 6c\)，那么（\quad）

\choices
    {\(\frac{1}{c} = \frac{1}{a} + \frac{1}{b}\)}
    {\(\frac{2}{c} = \frac{2}{a} + \frac{1}{b}\)}
    {\(\frac{1}{c} = \frac{2}{a} + \frac{2}{b}\)}
    {\(\frac{2}{c} = \frac{1}{a} + \frac{2}{b}\)}
\end{problem}

\begin{answer}
题面无正确选项
\end{answer}

\begin{solution}
令 \(3a=4b=6c=t\)，则 \(a=\frac{t}{3}\)、\(b=\frac{t}{4}\)、\(c=\frac{t}{6}\)，从而
\(\displaystyle \frac{1}{a}=\frac{3}{t}\)，\(\displaystyle \frac{1}{b}=\frac{4}{t}\)，\(\displaystyle \frac{1}{c}=\frac{6}{t}\). 逐项代入可得：第一项的右边为 \(\frac{7}{t}\) 而左边为 \(\frac{6}{t}\)；第二项的左、右两边分别为 \(\frac{12}{t}\)、\(\frac{10}{t}\)；第三项的左、右两边分别为 \(\frac{6}{t}\)、\(\frac{14}{t}\)；第四项的左、右两边分别为 \(\frac{12}{t}\)、\(\frac{11}{t}\). 四个等式均不成立，因此按原题面不存在正确选项；现有短答案不能由题面推出.
\end{solution}

\begin{problem}
同室四人各写一张贺年卡，先集中起来，然后每人从中拿一张别人送出的贺年卡，则四张贺年卡不同的分配方式有（\quad）

\choices
    {\(6\) 种}
    {\(9\) 种}
    {\(11\) 种}
    {\(23\) 种}
\end{problem}

\begin{answer}
B
\end{answer}

\begin{solution}
四张贺年卡的分配要求是每个人都不能拿到自己写的卡，即求 \(4\) 个元素的错排数. 用容斥原理，分配方式数为
\(\displaystyle 4!-\mathrm C_4^1 3!+\mathrm C_4^2 2!-\mathrm C_4^3 1!+\mathrm C_4^4 0!=24-24+12-4+1=9\)，故选 B.
\end{solution}

\begin{problem}
在正方体 \(A_{1}B_{1}C_{1}D_{1}-ABCD\) 中，\(M\)，\(N\) 分别为棱 \(A_{1}A\) 和棱 \(B_{1}B\) 的中点（如图）. 若 \(\theta\) 为直线 \(CM\) 与 \(D_{1}N\) 所成的角，则 \(\sin \theta\) 的值为（\quad）

\bitmapfigure[width=0.42\linewidth]{img/1993/national_liberal/q18_fig1.png}

\choices
    {\(\frac{1}{9}\)}
    {\(\frac{2}{3}\)}
    {\(\frac{2\sqrt{5}}{9}\)}
    {\(\frac{4\sqrt{5}}{9}\)}
\end{problem}

\begin{answer}
D
\end{answer}

\begin{solution}
取正方体棱长为 \(1\)，建立坐标系
\(A=(0,0,0)\)、\(B=(1,0,0)\)、\(C=(1,1,0)\)、\(D_1=(0,1,1)\)，并取 \(M=(0,0,\frac{1}{2})\)、\(N=(1,0,\frac{1}{2})\). 可取直线 \(CM\)、\(D_1N\) 的方向向量分别为 \(\bs u=(-2,-2,1)\)、\(\bs v=(2,-2,-1)\). 于是 \(\bs u\cdot\bs v=-1\)，且 \(|\bs u|=|\bs v|=3\). 两直线所成的角取锐角，故 \(\displaystyle \cos\theta=\frac{|\bs u\cdot\bs v|}{|\bs u||\bs v|}=\frac{1}{9}\)，从而 \(\displaystyle \sin\theta=\sqrt{1-\frac{1}{81}}=\frac{4\sqrt{5}}{9}\)，故选 D.
\end{solution}

\section{填空题}

\begin{problem}
抛物线 \(y^2 = 4x\) 的弦 \(AB\) 垂直于 \(x\) 轴，若 \(AB\) 的长为 \(4\sqrt{3}\)，则焦点到 \(AB\) 的距离为\fillinblank{}.
\end{problem}

\begin{answer}
2
\end{answer}

\begin{solution}
由抛物线方程 \(y^2=4x\) 可知焦点为 \(F=(1,0)\). 弦 \(AB\) 垂直于 \(x\) 轴，设其所在直线为 \(x=t\)，则弦端点的纵坐标为 \(\pm2\sqrt{t}\)，从而 \(AB=4\sqrt{t}=4\sqrt{3}\)，得 \(t=3\). 因此焦点到直线 \(AB\) 的距离为 \(|3-1|=2\).
\end{solution}

\begin{problem}
在半径为 \(30\mathrm{~m}\) 的圆形广场中央上空，设置一个照明光源，射向地面的光呈圆锥形，且其轴截面顶角为 \(120^{\circ}\). 若要光源恰好照亮整个广场，则其高度应为\fillinblank{} \(\mathrm{m}\)（精确到 \(0.1\mathrm{~m}\)）.
\end{problem}

\begin{answer}
17.3
\end{answer}

\begin{solution}
设光源到地面的高度为 \(h\). 圆锥轴截面的顶角为 \(120^\circ\)，所以半顶角为 \(60^\circ\). 光源、广场圆心和广场边缘构成直角三角形，故 \(\displaystyle\tan60^\circ=\frac{30}{h}\)，从而 \(h=\frac{30}{\sqrt{3}}=10\sqrt{3}\approx17.3205\). 精确到 \(0.1\mathrm{~m}\) 得 \(17.3\mathrm{~m}\).
\end{solution}

\begin{problem}
在 \(50\) 件产品中有 \(4\) 件是次品，从中任意抽出 \(5\) 件，至少有 \(3\) 件是次品的抽法共\fillinblank{}\(\text{种}\)（用数字作答）.
\end{problem}

\begin{answer}
4186
\end{answer}

\begin{solution}
抽出的 \(5\) 件中至少有 \(3\) 件次品，分为抽出 \(3\) 件次品或 \(4\) 件次品两种互斥情形. 前一种有 \(\mathrm C_4^3\mathrm C_{46}^2\) 种，后一种有 \(\mathrm C_4^4\mathrm C_{46}^1\) 种，因此总数为 \(\displaystyle\mathrm C_4^3\mathrm C_{46}^2+\mathrm C_4^4\mathrm C_{46}^1=4\cdot\frac{46\cdot45}{2}+46=4140+46=4186\).
\end{solution}

\begin{problem}
建造一个容积为 \(8\mathrm{~m}^{3}\)，深为 \(2\mathrm{~m}\) 的长方体无盖水池，如果池底和池壁的造价每平方米分别为 \(120\text{ 元}\) 和 \(80\text{ 元}\)，那么水池的最低总造价为\fillinblank{}\(\text{元}\).
\end{problem}

\begin{answer}
1760
\end{answer}

\begin{solution}
设池底的长、宽分别为 \(x\)，\(y>0\). 由容积和深度得 \(2xy=8\)，所以 \(xy=4\). 池底造价为 \(120xy=480\) 元；四面池壁的总面积为 \(2\cdot2x+2\cdot2y=4(x+y)\)，池壁造价为 \(320(x+y)\) 元. 因此总造价为 \(480+320(x+y)\). 由 \(x+y\geqslant2\sqrt{xy}=4\)，总造价不小于 \(480+320\cdot4=1760\) 元，且 \(x=y=2\) 时取等号，所以最低总造价为 \(1760\) 元.
\end{solution}

\begin{problem}
设 \(f(x) = 4^{x} - 2^{x + 1}\)，则 \(f^{-1}(0) =\)\fillinblank{}.
\end{problem}

\begin{answer}
1
\end{answer}

\begin{solution}
由 \(f^{-1}(0)\) 的定义，令 \(f(x)=0\). 于是 \(\displaystyle4^x-2^{x+1}=2^{2x}-2\cdot2^x=2^x(2^x-2)=0\). 由于 \(2^x>0\)，只能有 \(2^x=2\)，所以 \(x=1\)，即 \(f^{-1}(0)=1\).
\end{solution}

\begin{problem}
设 \(a > 1\)，则 \(\displaystyle\lim_{n \to \infty} \frac{1 - a^{n + 1}}{1 + a^{n + 1}} =\)\fillinblank{}.
\end{problem}

\begin{answer}
-1
\end{answer}

\begin{solution}
分子、分母同除以 \(a^{n+1}\)，得 \(\displaystyle\frac{1-a^{n+1}}{1+a^{n+1}}=\frac{a^{-(n+1)}-1}{a^{-(n+1)}+1}\). 因为 \(a>1\)，所以 \(a^{-(n+1)}\to0\)，从而极限为 \(\displaystyle\frac{0-1}{0+1}=-1\).
\end{solution}

\section{解答题}

\begin{problem}
解方程：\(\lg (x^2 + 4x - 26) - \lg (x - 3) = 1\).
\end{problem}

\begin{solution}
由对数真数条件，必须有 \(x-3>0\) 且 \(x^2+4x-26>0\). 前者给出 \(x>3\)，后者的根为 \(-2\pm\sqrt{30}\)，结合 \(x>3\) 得定义域为 \(x>-2+\sqrt{30}\). 在此条件下，原方程化为
\(\displaystyle\lg\frac{x^2+4x-26}{x-3}=1\)，所以 \(\frac{x^2+4x-26}{x-3}=10\). 整理得 \(x^2-6x+4=0\)，解得 \(x=3\pm\sqrt{5}\). 其中 \(3-\sqrt{5}<3\) 不满足定义域，\(3+\sqrt{5}>-2+\sqrt{30}\) 满足定义域. 因此原方程的解为 \(x=3+\sqrt{5}\).
\end{solution}

\begin{problem}
已知数列 \(\frac{8 \cdot 1}{1^2 \cdot 3^2}\)，\(\frac{8 \cdot 2}{3^2 \cdot 5^2}\)，\(\cdots\)，\(\frac{8n}{(2n-1)^2 (2n+1)^2}\)，\(\cdots\)，\(S_n\) 为其前 \(n\) 项和. 计算得 \(S_1 = \frac{8}{9}\)，\(S_2 = \frac{24}{25}\)，\(S_3 = \frac{48}{49}\)，\(S_4 = \frac{80}{81}\). 观察上述结果，推测出计算 \(S_n\) 的公式，并用数学归纳法加以证明.
\end{problem}

\begin{solution}
记数列的第 \(n\) 项为 \(a_n\). 因为
\(\displaystyle a_n=\frac{8n}{(2n-1)^2(2n+1)^2}=\frac{1}{(2n-1)^2}-\frac{1}{(2n+1)^2}\)，所以由前几项可猜想
\(\displaystyle S_n=1-\frac{1}{(2n+1)^2}=\frac{4n(n+1)}{(2n+1)^2}\).

当 \(n=1\) 时，\(\displaystyle S_1=\frac{8}{9}=1-\frac{1}{3^2}\)，公式成立. 假设 \(n=k\) 时公式成立，即 \(\displaystyle S_k=1-\frac{1}{(2k+1)^2}\)，则
\(\displaystyle S_{k+1}=S_k+a_{k+1}=1-\frac{1}{(2k+1)^2}+\frac{1}{(2k+1)^2}-\frac{1}{(2k+3)^2}=1-\frac{1}{(2k+3)^2}\)，这正是把 \(n\) 换成 \(k+1\) 后的公式. 因此由数学归纳法，对任意正整数 \(n\)，都有 \(\displaystyle S_n=1-\frac{1}{(2n+1)^2}=\frac{4n(n+1)}{(2n+1)^2}\).
\end{solution}

\begin{problem}
如图，\(A_{1}B_{1}C_{1}-ABC\) 是直三棱柱，过点 \(A_{1}\)，\(B\)，\(C_{1}\) 的平面和平面 \(ABC\) 的交线记作 \(l\).

\begin{enumerate}
\item 判定直线 \(A_{1}C_{1}\) 和 \(l\) 的位置关系，并加以证明；
\item 若 \(A_{1}A = 1\)，\(AB = 4\)，\(BC = 3\)，\(\angle ABC = 90^{\circ}\)，求顶点 \(A_{1}\) 到直线 \(l\) 的距离.
\end{enumerate}

\bitmapfigure[width=0.42\linewidth]{img/1993/national_liberal/q27_fig1.png}
\end{problem}

\begin{solution}
\begin{enumerate}
\item 设 \(\Pi\) 为过 \(A_1\)、\(B\)、\(C_1\) 的平面. 由直三棱柱的性质，平面 \(A_1B_1C_1\) 与平面 \(ABC\) 平行. 又 \(A_1C_1=\Pi\cap(A_1B_1C_1)\)，而 \(l=\Pi\cap ABC\)，所以同一平面与两平行平面的交线互相平行，即 \(A_1C_1\parallel l\).
\item 因为 \(A_1C_1\parallel AC\)，且 \(A_1C_1\parallel l\)，所以 \(l\parallel AC\). 设 \(E\in l\) 且 \(AE\perp l\). 由于 \(A_1A\perp\) 平面 \(ABC\)，直线 \(AE\) 是斜线 \(A_1E\) 在平面 \(ABC\) 上的射影；由三垂线定理的逆定理，得 \(A_1E\perp l\)，所以 \(A_1E\) 是所求距离. 在直角三角形 \(A_1AE\) 中，
\(A_1E^2=A_1A^2+AE^2\).

在直角三角形 \(ABC\) 中，\(AC=\sqrt{AB^2+BC^2}=\sqrt{4^2+3^2}=5\)，面积为 \(\frac{1}{2}\cdot4\cdot3=6\). 因 \(l\parallel AC\) 且 \(B\in l\)，所以 \(AE\) 等于点 \(B\) 到直线 \(AC\) 的距离. 设该距离为 \(BD\)，则 \(\frac{1}{2}\cdot AC\cdot BD=6\)，从而 \(BD=AE=\frac{12}{5}\). 因此
\(\displaystyle A_1E=\sqrt{1^2+\left(\frac{12}{5}\right)^2}=\frac{13}{5}\).
\end{enumerate}
\end{solution}

\begin{problem}
在面积为 \(1\) 的 \(\triangle PMN\) 中，\(\tan \angle PMN = \frac{1}{2}\)，\(\tan \angle MNP = -2\)，建立适当的坐标系，求出以 \(M\)，\(N\) 为焦点且过点 \(P\) 的椭圆方程.
\end{problem}

\begin{solution}
以 \(MN\) 所在直线为 \(x\) 轴，以 \(MN\) 的垂直平分线为 \(y\) 轴，设 \(M=(-c,0)\)、\(N=(c,0)\)、\(P=(x,y)\)，并取 \(y>0\). 由 \(\tan\angle PMN=\frac{1}{2}\) 得 \(\frac{y}{x+c}=\frac{1}{2}\)；由于 \(\tan\angle MNP=-2\) 表示点 \(P\) 在 \(N\) 的右上方，得 \(\frac{y}{x-c}=2\). 解方程组
\(\displaystyle y=\frac{x+c}{2}\)，\(\displaystyle y=2(x-c)\)，解得 \(x=\frac{5c}{3}\)、\(y=\frac{4c}{3}\).

三角形 \(PMN\) 的底 \(MN=2c\)，高为 \(y=\frac{4c}{3}\)，由面积为 \(1\) 得 \(\displaystyle\frac{1}{2}\cdot2c\cdot\frac{4c}{3}=1\)，所以 \(c^2=\frac{3}{4}\).

设所求椭圆的长半轴为 \(a\)，短半轴为 \(b\)，则 \(a^2=b^2+c^2\)，且 \(P\) 在椭圆上. 令 \(u=\frac{a^2}{c^2}\)，代入 \(\displaystyle\frac{x^2}{a^2}+\frac{y^2}{b^2}=1\) 得
\(\displaystyle\frac{25}{9u}+\frac{16}{9(u-1)}=1\)，即 \(9u^2-50u+25=0\). 解得 \(u=5\) 或 \(u=\frac{5}{9}\)，因椭圆必须满足 \(a^2>c^2\)，舍去 \(u=\frac{5}{9}\)，故 \(a^2=5c^2=\frac{15}{4}\)，\(b^2=a^2-c^2=3\). 所以椭圆方程为
\(\displaystyle\frac{x^2}{15/4}+\frac{y^2}{3}=1\).
\end{solution}

\begin{problem}
设复数 \(z = \cos \theta + \i\sin \theta\)（\(0 < \theta < \pi\)），\(\omega = \frac{1 - \left( \overline{z} \right)^4}{1 + z^4}\)，已知 \(|\omega| = \frac{\sqrt{3}}{3}\)，\(\arg \omega = \frac{\pi}{2}\)，求 \(\theta\).
\end{problem}

\begin{solution}
由 \(z=\cos\theta+\i\sin\theta\) 得 \(z=\e^{\i\theta}\)、\(\overline{z}=\e^{-\i\theta}\). 原式有意义要求 \(1+z^4\ne0\)，即 \(\cos2\theta\ne0\). 因而
\(\displaystyle\omega=\frac{1-\e^{-4\i\theta}}{1+\e^{4\i\theta}}=\frac{2\i\e^{-2\i\theta}\sin2\theta}{2\e^{2\i\theta}\cos2\theta}=\i\e^{-4\i\theta}\tan2\theta=\tan2\theta\left(\sin4\theta+\i\cos4\theta\right)\).

由 \(|\omega|=\frac{\sqrt{3}}{3}>0\)，得 \(\tan2\theta\ne0\). 又 \(\arg\omega=\frac{\pi}{2}\) 且 \(|\omega|>0\)，所以 \(\omega\) 是正纯虚数，实部必须为零. 因此
\(\tan2\theta\sin4\theta=0\)，从而 \(\sin4\theta=0\). 结合 \(0<\theta<\pi\)，候选值只有 \(\theta=\frac{\pi}{4}\)、\(\frac{\pi}{2}\)、\(\frac{3\pi}{4}\).

当 \(\theta=\frac{\pi}{4}\) 或 \(\frac{3\pi}{4}\) 时，\(\cos2\theta=0\)，即 \(1+z^4=0\)，原式无意义；当 \(\theta=\frac{\pi}{2}\) 时，\(\tan2\theta=0\)，从而 \(\omega=0\)，与 \(|\omega|=\frac{\sqrt{3}}{3}\) 矛盾. 因此按原始题面不存在满足全部条件的 \(\theta\)，现题条件彼此不相容.
\end{solution}
